\section{Consequences and the Fifth Target}
\label{sec:consequences}
\begin{corollary}\label{cor:three}
Binary linear $[44,10,\ge18]$, $[44,11,\ge17]$, and $[45,11,\ge18]$ codes do not exist.
\end{corollary}
\begin{proof}
Puncturing a $[44,10,\ge18]$ code gives a $[43,10,\ge17]$ code, contradicting Theorem~\ref{thm:main}. Shortening a $[44,11,\ge17]$ code on any coordinate gives a length-43 code of dimension at least ten and distance at least 17; a ten-dimensional subcode gives the same contradiction. Finally, puncturing a $[45,11,\ge18]$ code gives the excluded $[44,11,\ge17]$ parameters.
\end{proof}

For the fifth target, the exact-distance convention in E3 must be respected. The following elementary lemma supplies the connection.
\begin{lemma}[Distance normalization]\label{lem:normalize}
For $k\ge1$ and $d\ge1$, existence of a binary linear $[n,k,\ge d]$ code implies existence of a binary linear $[n,k,d]$ code.
\end{lemma}
\begin{proof}
Let the original minimum distance be $D\ge d$, and choose a word of weight $D$. Puncture $D-d$ positions in its support. Since $D-d<D$, puncturing is injective. Every nonzero word retains weight at least $D-(D-d)=d$, and the chosen word attains $d$. Append $D-d$ zero coordinates to restore length $n$ without changing dimension or distance.
\end{proof}

\begin{theorem}\label{thm:42}
There is no binary linear $[42,16,\ge13]$ code.
\end{theorem}
\begin{proof}
Parity extension of a hypothetical code gives $[43,16,\ge14]$. Shortening one coordinate loses at most one dimension. Choose a 15-dimensional subcode, shorten another coordinate, and choose a 14-dimensional subcode. The result has parameters $[41,14,\ge14]$. Lemma~\ref{lem:normalize} gives exact $[41,14,14]$, contradicting E3, namely Proposition~2 of \citet{bouyuklieva2020}.
\end{proof}
The construction in this proof permits zero columns. This is consistent with both our original domain and the binary linear code domain of E3. In particular, the proof does not assume that a minimum-weight witness survives arbitrary shortening; normalization is performed after the shortening steps.

\begin{corollary}\label{cor:optima}
The optimal distances satisfy
\begin{align*}
 \dopt(43,10)&=16,&\dopt(44,10)&=17,\\
 \dopt(44,11)&=16,&\dopt(42,16)&=12,
\end{align*}
and $16\le\dopt(45,11)\le17$.
\end{corollary}
\begin{proof}
Theorems~\ref{thm:main} and~\ref{thm:42} and Corollary~\ref{cor:three} give the upper bounds. Appendix~\ref{app:lower} specifies full-rank generators of $[43,10,16]$, $[44,10,17]$, $[44,11,16]$, and $[42,16,12]$, together with their fully enumerated weight spectra. Appending a zero coordinate to the $[44,11,16]$ code gives $[45,11,16]$.
\end{proof}
